\chapter{Sprint 57 --- Cụm D (bù hàng, font, CardHeader) + Onboarding cửa hàng (2026-08-25 $\to$ 2026-09-04)}

\emph{Chapter này gộp toàn bộ sprint D: bốn cụm giao diện D1--D3 chạy từ 25/08 đến 04/09,
và issue nghiệp vụ \texttt{WJ-FRANCHISE-003} đóng ngày 04/09. Số đo chi tiết từng cụm nằm ở
\texttt{docs/d\{,b,c,d,e,f\}-acceptance-matrix.md}, \texttt{docs/d2-acceptance-matrix.md},
\texttt{docs/d3-review-matrix.md} và \texttt{docs/wj-franchise-003-acceptance-matrix.md}.}

\section{Bối cảnh}

Sau khi Sprint 56 đóng batch C1--C10, BA audit UAT mở lứa issue mới. Phiên 25/08 phân thành
sáu cụm \texttt{D1--D6} cộng năm cụm refactor \texttt{R1--R5}
(\texttt{docs/next-session-clusters-D.md}). Sprint này chạy hết D1, D2, D3 và chèn thêm một
issue nghiệp vụ Severity High mà BA mở ngày 02/09.

\section{D1 --- Bù hàng: bốn issue nhãn và bộ lọc (25/08)}

Bốn issue \texttt{UAT-BH-001/003/005/006}, module \texttt{wujia\_portal\_return}. Gốc rễ
chung: nhãn hiển thị và nhãn trong ô lọc được viết ở hai chỗ khác nhau nên trôi khỏi nhau ---
danh sách chọn có \emph{hai dòng cùng ghi} ``Đang xử lý'', còn ``Đang bù một phần'' hiện được
trên thẻ nhưng không lọc được vì nó không phải một trạng thái lưu trong bảng mà là kết hợp
của \emph{đang xử lý} và \emph{đã giao một phần}. Cách sửa là gom về \textbf{một danh sách
nhãn duy nhất} cho cả hiển thị lẫn bộ lọc, và định nghĩa hai nhãn chồng lấn thành hai tập
rời nhau. Kiểm chứng bằng một tính chất chứ không bằng mắt: dựng dữ liệu đủ tám trạng thái
rồi cộng số kết quả của tám lần lọc --- phải đúng bằng tổng danh sách khi không lọc, tức là
các nhãn vừa không trùng vừa không bỏ sót.

\section{D2 --- Font Portal về Inter toàn bộ (26/08)}

Issue \texttt{WJ-PORTAL-UI-002}. BA nêu bốn màn còn hiện font Inter Tight. Đo trên UAT cho
thấy đó không phải lỗi của bốn màn, mà là \textbf{một luật CSS đặt hụt phạm vi}: luật ép font
chỉ có hiệu lực bên trong khung nội dung chính, trong khi các màn mobile và khối danh sách
Khảo sát lại dựng bên ngoài khung đó. Nới phạm vi luật ra toàn trang sửa một lần cho tất cả,
và máy quét lộ thêm ba màn nữa cùng bệnh mà BA chưa thấy. Kèm theo: khai font dự phòng cho
chữ Thái và chữ Trung, vì bộ Inter tự cài chỉ có Latinh và tiếng Việt.

\section{D3 --- CardHeader \texttt{CMP-CH-001}: sáu cụm migrate + một phiên review (27/08 -- 04/09)}

Cụm lớn nhất sprint. Kiểm kê \texttt{docs/d3-cardheader-inventory.md} lọc 218 ứng viên xuống
\textbf{103 chỗ gọi thật}, phân loại \textbf{theo tổ tiên DOM chứ không theo tên class} (kế
thừa bài học C8), trong đó có 31 ``giả-heading'' là thẻ \texttt{<p>}/\texttt{<span>} đóng vai
tiêu đề --- đúng bệnh BA nêu.

\begin{itemize}
\item \textbf{D3a} 27/08 --- dựng \texttt{wj\_card\_header} + 10 chỗ gọi mẫu.
\item \textbf{D3b} 28/08 --- 26 chỗ gọi, phủ 36/103.
\item \textbf{D3c} 02/09 --- 11 chỗ gọi, phủ 47/103.
\item \textbf{D3d} 04/09 --- 14 chỗ gọi ở màn thi, phủ 61/103.
\item \textbf{D3e} 04/09 --- 24 chỗ gọi ở đổi trả và lịch sử, phủ 85/103.
\item \textbf{D3f} + review 04/09 --- 10 chỗ gọi cuối, \textbf{95/105}, rồi một phiên đo lại
  toàn cụm bằng ảnh chụp trên chính UAT.
\end{itemize}

Bốn bài học đắt nhất của cụm này, đáng nhớ hơn cả con số:

\begin{enumerate}
\item \textbf{Số đo Pass hết mà bố cục vẫn vỡ.} Ở D3c, sau migrate thì mọi chỉ số cỡ chữ,
  dòng cao, màu, khoảng cách đều đạt --- nhưng nhãn trạng thái \textbf{trôi ra mép phải cột}
  vì vùng dẫn của component là \texttt{flex:1 1 auto} nên nở hết bề rộng. Chỉ ảnh chụp mới
  bắt được. Từ đó chốt: cụm giao diện nào cũng phải có một lượt review bằng ảnh.
\item \textbf{\texttt{:not()} mang độ đặc hiệu của tham số.} Ở D3b, tiêu đề trong
  \texttt{.card > .card-body} ra màu \texttt{\#212529} thay vì \texttt{\#111827}. Giả thuyết
  ``thiếu token'' \textbf{sai}; thủ phạm truy bằng CDP là một luật của theme Vuexy đạt
  \texttt{(0,4,1)} nhờ \texttt{:not()}, đè luật của component ở \texttt{(0,1,0)}.
\item \textbf{Đo rồi mới thêm luật --- và kết luận là không thêm gì.} Sáu lớp thân thẻ
  ``ứng viên'' bị nghi cộng chồng khoảng cách; đo thật thì \emph{không lớp nào} cộng chồng,
  thêm luật mù sẽ bóp nhịp xuống dưới chuẩn BA.
\item \textbf{Harness sai còn nguy hơn không đo.} Lượt đo UAT đầu của D3b báo chín cảnh báo
  thì chỉ ba là lỗi thật: ba cái do đường dẫn không tồn tại nên Odoo trả trang danh sách và
  máy đo nhầm trang, ba cái do kỳ vọng viết quá rộng.
\end{enumerate}

Phiên review 04/09 (\texttt{0d8366d}) đo \textbf{quan hệ} thay vì đo tuyệt đối: 31 đường dẫn
$\times$ 4 khổ màn = 124 lượt, 174 tiêu đề trong thẻ --- 0 lỗi kịch bản, 0 tràn ngang, 0 nhóm
cỡ chữ trôi chưa giải trình. Sửa ba chỗ, trong đó có một lỗi tương phản \textbf{có sẵn}
(3,74 lên 5,42). Xoá 17 lớp CSS chết. Bài học D3c của phiên review: chứng minh bản cài đặt ăn
thật phải bằng \emph{giá trị chỉ có sau khi sửa}, không phải bằng việc nhìn thấy bảng sạch cờ.

Issue \texttt{UI-CARDHEADER-001} \textbf{chưa đóng}: BA đã tự chuyển sang
\texttt{Need Clarification} nên Dev không ghi sheet được, và bản thân nó vẫn ở 95/105 vì chín
chỗ còn lại đều là câu hỏi treo chờ BA.

\section{\texttt{WJ-FRANCHISE-003} --- Onboarding cửa hàng một luồng (04/09)}

Issue Severity High, BA mở 02/09. Mở một cửa hàng mới đang phải đi qua ba màn rời: tạo hồ sơ
cửa hàng, tự tìm rồi gắn đối tác, sang phần Người dùng tạo tài khoản và phải tự biết chọn
đúng loại tài khoản (có khi phải bật chế độ kỹ thuật), rồi quay lại gán thành viên. Hậu quả:
dễ đẻ đối tác và tài khoản trùng, gắn nhầm đối tác cửa hàng thành hồ sơ cá nhân, bỏ sót thành
viên.

Commit \texttt{8a4932e}, module \texttt{wujia\_franchise} \texttt{19.0.3.1.1} $\to$
\texttt{19.0.4.0.0}. \textbf{Không bảng mới, không đổi cấu trúc dữ liệu, không migration,
Portal không đụng một dòng.}

\begin{itemize}
\item \texttt{wujia.franchise.onboarding.wizard} (transient) + dòng con
  \texttt{...member.line}: một màn gộp bốn việc, vào được từ menu \emph{Vận hành $\to$
  Onboarding cửa hàng} và từ hai nút trên hồ sơ cửa hàng.
\item \texttt{status} đổi mặc định \texttt{active} $\to$ \texttt{draft}; bản ghi cũ giữ
  nguyên.
\item \texttt{\_assert\_ready\_to\_activate()} gọi từ \texttt{action\_set\_active()}: phải có
  đối tác và ít nhất một chủ tiệm hiệu lực. Đặt ở \textbf{method hành động}, không phải
  \texttt{@api.constrains} --- bài học L12: constrains sẽ khoá cứng cửa hàng cũ, sửa ghi chú
  cũng không nổi.
\item Dò trùng đối tác theo tên / 8 số cuối điện thoại / thư điện tử / mã số thuế, chặn tới
  khi HQ tick xác nhận.
\item Ba lớp chặn cho tài khoản sẵn có: nhóm nội bộ, đã lưu trữ, đã là thành viên --- báo lỗi
  rõ và \textbf{tuyệt đối không tự sửa quyền của ai}.
\item Quyền: \texttt{\_assert\_hq()} kiểm nhóm rồi mới \texttt{sudo()}, thay vì cấp quyền
  quản trị người dùng cho HQ (rộng hơn nhiều so với điều BA cần).
\item Tái dùng: \texttt{res.partner.\_wujia\_franchise\_mapping()} của cụm C1 cho cảnh báo
  mapping, \texttt{main\_owner\_member\_id} sẵn có cho điều kiện kích hoạt, và tách
  \texttt{ROLE\_SELECTION} ra hằng dùng chung thay vì chép lại danh sách vai trò.
\end{itemize}

\paragraph{Lý do / Trade-off.} Hai quyết định đáng ghi. Thứ nhất, \textbf{hồ sơ cá nhân của
tài khoản Portal đứng độc lập, không làm con của đối tác cửa hàng}: gắn cha thì danh bạ gọn
hơn, nhưng khi đó ``đối tác thương mại'' của người dùng chính là cửa hàng, và bộ quy tắc suy
ra cửa hàng từ đối tác (dựng ở cụm C1) sẽ tự gán cửa hàng cho mọi chứng từ mang tên cá nhân
đó --- lẫn công nợ. Thứ hai, \textbf{cửa hàng Nháp mới chỉ bị chặn ở Portal}: lớp chặn đơn
hàng và hoá đơn nằm trong code vừa nghiệm thu ở C1/C2, đụng vào là phải chạy lại toàn bộ hồi
quy cụm ấy, nên ghi thành LIMIT tường minh cho BA quyết.

\paragraph{Chứng minh.} 18 phép thử mới, \textbf{bảy lớp kiểm tra đều được chứng minh có tác
dụng bằng mutation} --- tạm gỡ từng lớp ra thì đúng phép thử tương ứng đỏ. Đáng nói là lượt
mutation đầu tiên phát hiện lớp chặn ``tài khoản nội bộ'' \emph{không được chứng minh}: gỡ nó
ra mà không phép thử nào đỏ, vì lớp kế tiếp bắt hộ. Phải siết phép thử để nó kiểm đúng nội
dung thông báo thì mutation mới đỏ --- đúng bài học L8 rằng khẳng định lỏng là khẳng định
rỗng.

\paragraph{Hồi quy chứng minh bằng đối chứng.} Chạy 13 module: bản có thay đổi ra
\texttt{0 failed, 7 error / 299}; chạy \emph{đúng lệnh đó} trên bản chưa thay đổi (stash) ra
\texttt{0 failed, 7 error / 282}. Bảy lỗi giống hệt $\Rightarrow$ có sẵn, do
\texttt{wujia\_portal\_inspection} chưa cài trên DB kiểm thử. Không có lượt đối chứng này thì
bảy lỗi đó rất dễ bị quy oan cho đợt sửa.

\paragraph{Đo trước khi bật cổng.} Đọc chỉ-đọc trên UAT qua XML-RPC: 3 cửa hàng, cả 3 đang
hoạt động, 0 cửa hàng thiếu đối tác, \textbf{1 cửa hàng \texttt{HN-02} thiếu chủ tiệm hiệu
lực}. Vì điều kiện đặt ở nút nên \texttt{HN-02} vẫn chạy bình thường.

\paragraph{Ghi lại cho BA.} Spec để mở nhiều chi tiết kỹ thuật, nên toàn bộ \textbf{12 điểm
Dev phải tự quyết} được ghi thành một bản riêng
(\texttt{docs/wj-franchise-003-dev-decisions.md}), mỗi mục nêu \emph{spec nói gì $\cdot$ Dev
chọn gì $\cdot$ vì sao $\cdot$ đổi thì mất gì}. Trong đó có một phát hiện ngoài phạm vi:
màn phiếu giám sát cửa hàng cũng đang tự sinh tài khoản Portal bằng cách so khớp \emph{theo
tên} rồi đặt tên đăng nhập ngẫu nhiên --- đúng nguồn sinh tài khoản trùng mà issue này muốn
dẹp, nhưng nằm ở luồng khác nên đề nghị BA mở issue riêng.

\paragraph{Dịch màn hình mới sang tiếng Việt --- ba bẫy im lặng của i18n Odoo 19.}
Phát sinh \emph{sau} khi cài lên UAT ngày 04/09: toàn bộ màn onboarding hiện tiếng Anh giữa một
backend tiếng Việt. Người dùng màn này là nhân sự HQ người Việt nên đây không phải chuyện thẩm
mỹ. Dev chọn dịch bằng file ngôn ngữ \textbf{trong mã nguồn} (\texttt{i18n/vi\_VN.po}) chứ không
dịch tay trên giao diện, vì bản dịch tiếng Việt đang có của module nằm \emph{trong database} ---
dựng lại DB là mất sạch, và bản chạy thật sau này sẽ phải dịch lại từ đầu. Ba bẫy gặp phải, cả
ba đều \textbf{không lỗi, không cảnh báo, chỉ là nhãn vẫn tiếng Anh}:

\begin{enumerate}
\item \textbf{Odoo ghép \texttt{.po} với \texttt{.pot} cùng thư mục lúc nạp.} Chuỗi có trong
  \texttt{.po} mà \textbf{thiếu trong \texttt{.pot}} thì bị \textbf{bỏ qua im lặng}. Sinh
  \texttt{.po} thì phải sinh \texttt{.pot} cùng lúc, không được sinh mỗi một file. Đây là
  \textbf{tái phát của L8(3)} đã ghi từ Sprint 44 --- lần đó biểu hiện là msgid mới thành
  obsolete, lần này là chuỗi mới không bao giờ được nạp; cùng một gốc.
\item \textbf{Dạng tham chiếu menu đã đổi.} \texttt{\#: menu:<module>.<xmlid>} \textbf{không
  còn hợp lệ}, phải viết \texttt{\#: model:ir.ui.menu,name:<module>.<xmlid>}. Bằng chứng còn
  nguyên trong chính \texttt{custom/wujia\_franchise/i18n/vi\_VN.po}: hai dạng cùng tồn tại
  cạnh nhau, và chỉ dạng \texttt{model:ir.ui.menu,name:} là ăn --- đọc menu \texttt{id 463}
  trên UAT ở \texttt{lang=vi\_VN} ra \emph{``Onboarding cửa hàng''}, còn các menu khai theo
  dạng cũ vẫn ra tiếng Anh.
\item \textbf{Chuỗi trong code cần \emph{chú thích đánh dấu}, không phải chỉ cần tham chiếu.}
  Sau bản \texttt{19.0.4.0.1}, tiêu đề hộp thoại và \emph{mọi} \texttt{UserError} vẫn tiếng
  Anh dù entry đã có tham chiếu \texttt{\#: code:addons/...}. Odoo 19 chỉ nhận bản dịch code
  khi entry mang dòng chú thích \texttt{\#. odoo-python} --- thiếu dòng đó thì tham chiếu
  \texttt{code:} là vô nghĩa. Phải thêm dấu này cho cả 27 entry code và cấp cho
  \emph{``Store onboarding''} một tham chiếu code riêng (nó vừa là tên menu vừa là tiêu đề hộp
  thoại dựng từ Python). Sửa ở \texttt{7568711}, module lên \texttt{19.0.4.0.2}.
\end{enumerate}

Bài học chung của ba bẫy: \textbf{i18n Odoo hỏng thì hỏng không tiếng động} --- không có cách
nào biết ngoài việc \emph{mở đúng màn ở đúng ngôn ngữ mà nhìn}, hoặc đọc chỉ-đọc từng nhãn qua
XML-RPC với \texttt{context=\{'lang': 'vi\_VN'\}}. Không được coi ``RC=0, cài sạch'' là bằng
chứng đã dịch xong.

Việc kèm theo: bỏ đánh số ``1. / 2. / 3.'' ở tiêu đề ba khối --- khi mở ở chế độ \emph{thêm
người dùng} thì khối 1 và 2 bị ẩn, chỉ còn mỗi nhãn ``3.'' đứng trơ. Tiêu đề hộp thoại nay đổi
theo chế độ: \emph{Onboarding cửa hàng} khi tạo mới, \emph{Thêm người dùng cửa hàng} khi bổ sung
người dùng.

\section{Trạng thái cuối sprint}

\begin{itemize}
\item Đã handoff cho BA: \texttt{UAT-BH-001/003/005/006} (D1),
  \texttt{WJ-PORTAL-UI-002/003/004} (D2), \texttt{WJ-FRANCHISE-003}.
\item Còn treo: \texttt{UI-CARDHEADER-001} ở \texttt{Need Clarification} --- cần gửi BA một
  lượt gộp năm câu hỏi để gỡ.
\item \textbf{Deploy UAT --- một phần.} \texttt{wujia\_franchise 19.0.4.0.1} đã lên UAT
  (đo chỉ-đọc 04/09: module \texttt{installed}, menu \emph{Onboarding cửa hàng} và toàn bộ
  nhãn trường của wizard đọc ở \texttt{lang=vi\_VN} đều ra tiếng Việt). Bản
  \textbf{\texttt{19.0.4.0.2}} (\texttt{7568711}, dấu \texttt{\#. odoo-python}) \textbf{chưa
  lên} --- đo trên UAT thấy hai tiêu đề hộp thoại dựng từ Python vẫn trả về
  \emph{``Store onboarding''} / \emph{``Add store users''} kể cả ở \texttt{lang=vi\_VN}. Lệnh
  \texttt{-u wujia\_franchise}, không cần đổi số hiệu bộ giao diện.
\item Đã gửi BA một văn bản gộp năm câu hỏi treo của \texttt{UI-CARDHEADER-001}
  (\texttt{docs/ba-questions-ui-cardheader-001.md}), kèm hai đề nghị mở issue riêng: chỗ tự
  sinh tài khoản theo tên ở màn phiếu giám sát, và cửa hàng \texttt{HN-02} thiếu chủ tiệm.
\end{itemize}
